% tikz-hops v1.0 2026/10/09
% Copyright (C) 2026 Yusuke Izawa
% This work may be distributed and/or modified under the conditions of the
% LaTeX Project Public License, version 1.3c or later. Maintained by Yusuke Izawa.
\documentclass[a4paper]{article}
\usepackage[margin=25mm]{geometry}
\usepackage{tikz-hops}
\usetikzlibrary{arrows.meta}
\usepackage{fancyvrb,booktabs,array,hyperref}
\newcommand\key[1]{\texttt{#1}}
\newcommand\pkg[1]{\textsf{#1}}
\newenvironment{example}
  {\VerbatimOut{\jobname.ex}}
  {\endVerbatimOut\par\medskip\noindent
   \VerbatimInput[frame=single,fontsize=\small]{\jobname.ex}\par\smallskip
   \noindent\input{\jobname.ex}\par\medskip}
\title{The \pkg{tikz-hops} package\\\large Line hops at path crossings for TikZ}
\author{Yusuke Izawa}
\date{v1.0, 2026/10/09}
\begin{document}
\maketitle

\section{What it does}
A path drawn with \key{hop over} cuts a gap at every crossing with a saved
path and fills the gap with a small arc.
The package is a TikZ library built on \pkg{spath3}.
It needs TeX Live 2022 or later and \pkg{spath3} 2.8 or later.

\begin{Verbatim}[frame=single,fontsize=\small]
\usepackage{tikz-hops}      % or \usetikzlibrary{hops}
\end{Verbatim}
\noindent This is the only line to add to a preamble. It loads
\pkg{tikz} and the library \texttt{hops}.

\section{Usage}
Save the path to hop over with \key{spath/save}, then name it in
\key{hop over}.
The saved path may be drawn or not.

\begin{example}
\begin{tikzpicture}
  \draw[thick, spath/save=wire] (1,-1) -- (1,1);
  \draw[->, hop over=wire] (0,0) -- (3,0);
\end{tikzpicture}
\end{example}

\key{hop over} takes a comma list.
All crossings are processed in one pass, so the result does not depend on the
order of the names.

\begin{example}
\begin{tikzpicture}
  \path[spath/save=a] (1,-1) -- (1,1);
  \path[spath/save=b] (2,-1) -- (2,1);
  \draw[gray] (1,-1) -- (1,1) (2,-1) -- (2,1);
  \draw[hop over={a,b}] (0,0) -- (3,0);
\end{tikzpicture}
\end{example}

The macro \verb|\drawhop| is a shorthand.
\begin{Verbatim}[frame=single,fontsize=\small]
\drawhop[<options>]{<names>}{<path>}
=  \draw[hop over={<names>}, <options>] <path>;
\end{Verbatim}
It keeps the calling convention of the first version of this code.

\begin{example}
\begin{tikzpicture}
  \path[spath/save=wire] (1,-1) -- (1,1);
  \drawhop[red]{wire}{(0,0) -- (2,0)}
\end{tikzpicture}
\end{example}

\section{Keys}
All keys live in \texttt{/tikz/hops/}.
Write \verb|hops/gap=2mm| inside any \verb|\draw| option list, or group keys
with \verb|hops={gap=2mm, shape=square}|.
Keys are read when the path ends, so their position in the option list does not
matter.

\begin{center}
\begin{tabular}{@{}>{\ttfamily}l>{\ttfamily}l p{0.45\linewidth}@{}}
\toprule
Key & Default & Meaning\\
\midrule
hop over=<names> & & Comma list of saved paths to hop over. Starts the hops.\\
hops/gap=<dim> & 1.6mm & Chord length cut at each crossing. Smaller than the distance between crossings.\\
hops/height=<dim> & gap/2 & Apex height of the hop. The default is a semicircle.\\
hops/shape=arc|square|gap & arc & \texttt{gap} cuts without inserting anything.\\
hops/side=upright|left|right & upright & Side of the bump, see below.\\
\bottomrule
\end{tabular}
\end{center}

Lengths are canvas lengths. They are not changed by \verb|scale| or other
coordinate transformations, like \key{line width}.

\begin{example}
\begin{tikzpicture}
  \path[spath/save=v] (1,-1) -- (1,1);
  \draw[gray] (1,-1) -- (1,1);
  \draw[hop over=v, hops/gap=4mm]                    (0, .8) -- (2, .8);
  \draw[hop over=v, hops/gap=4mm, hops/height=1mm]   (0, .4) -- (2, .4);
  \draw[hop over=v, hops={gap=4mm, shape=square}]    (0,  0) -- (2,  0);
  \draw[hop over=v, hops={gap=4mm, shape=gap}]       (0,-.4) -- (2,-.4);
\end{tikzpicture}
\end{example}

\subsection{Side}
\key{left} and \key{right} are relative to the direction of travel.
\key{upright} keeps the bump upward on horizontal segments. On vertical
segments it bends to the left of the direction of travel.

\begin{example}
\begin{tikzpicture}
  \path[spath/save=v] (1,-1) -- (1,1);
  \draw[gray] (1,-1) -- (1,1);
  \foreach \s [count=\i] in {upright,left,right}
    \draw[hop over=v, hops/side=\s] (0,1-\i*.5) -- (2,1-\i*.5);
\end{tikzpicture}
\end{example}

\section{Limits}
\begin{itemize}
\item The drawn path is one connected path. A path with several subpaths gets
  a gap between them.
\item A crossing closer to a path end than \key{gap}/2, and two crossings
  closer than \key{gap}, are not handled. The line is cut without a hop.
  Reduce \key{gap} or move the crossing.
\item A crossing at a corner of the path gives a diagonal chord.
\item The gap is measured along the path and the hop is placed on the chord.
  On strongly curved paths use a small \key{gap}.
\item Self-crossings are not supported.
\item A name that is not a saved path gives the warning
  \texttt{Package tikz-hops Warning} and the path is drawn without hops.
\item Arrow tips come from the \verb|\draw| options and appear only at the path ends.
\end{itemize}

\section{Related packages}
The \texttt{knots} library of \pkg{spath3} cuts over/under gaps without hops.
\pkg{circuitikz} has a \texttt{jump crossing} component for circuit wires.

\section{Deprecated}
\verb|\hopgap| is the former default gap. Redefine it to change the default
\key{hops/gap}. It is removed in 2.0.

\section{License}
LaTeX Project Public License 1.3c or later.
\end{document}
